% SEGMENT OLP-0486-S01
% Part: applied-modal-logics
% Chapter: epistemic-logic
% Section: truth-at-w

\documentclass[../../../include/open-logic-section]{subfiles}

\begin{document}

\olfileid{aml}{el}{trw}

\olsection{ஓர் உலகில் மெய்மை}

இயல்பான வாய்ப்புநிலைத் தருக்கத்தைப் போலவே, ஒவ்வோர் அறிவுநிலை மாதிரியும்
அதிலுள்ள எந்த உலகங்களில் எந்த !!{formula}s மெய்யாகின்றன என்பதைத்
தீர்மானிக்கிறது. தொடர்புமுறைப் பொருண்மையியலின் அடிப்படைக் கருத்துக்கு,
``மாதிரி~$\mModel{M}$ உலகம்~$w$ இல் !!{formula}~$!A$ ஐ மெய்யாக்குகிறது''
என்ற அதே குறியீட்டைப் பயன்படுத்துகிறோம். இத்தொடர்பு தொகுமுறையில்
வரையறுக்கப்படுகிறது; வாய்ப்புநிலையற்ற எல்லாச் செயற்குறிகளுக்கும் அது இயல்பான
வாய்ப்புநிலை நேர்வுடன் ஒன்றுபட்டது.

\begin{defn}\ollabel{defn:mmodels}
  ஒரு~$\mModel M$ இல் \emph{$w$ இல் !!a{formula}~$!A$ இன் மெய்மை},
  குறியீட்டில்: $\mSat{M}{!A}[w]$, பின்வருமாறு தொகுமுறையில்
  வரையறுக்கப்படுகிறது:
  \begin{enumerate}
  \tagitem{prvFalse}{\indcase{!A}{\lfalse}{$\mSat{M}{\lfalse}[w]$ ஒருபோதும் மெய்யன்று}.}{}
  \tagitem{prvTrue}{\indcase{!A}{\ltrue}{$\mSat{M}{\ltrue}[w]$ எப்போதும் மெய்}.}{}
  \item $\mSat{M}{p}[w]$ எனில், எனில் மட்டுமே, $w \in V(p)$
  \tagitem{prvNot}{\indcase{!A}{\lnot !B}{$\mSat{M}{\indfrm}[w]$ எனில்,
    எனில் மட்டுமே, $\mSat/{M}{!B}[w]$}.}{}
  \tagitem{prvAnd}{\indcase{!A}{(!B \land !C)}{$\mSat{M}{\indfrm}[w]$ எனில்,
    எனில் மட்டுமே, $\mSat{M}{!B}[w]$ மற்றும் $\mSat{M}{!C}[w]$}.}{}
  \tagitem{prvOr}{\indcase{!A}{(!B \lor !C)}{$\mSat{M}{\indfrm}[w]$ எனில்,
    எனில் மட்டுமே, $\mSat{M}{!B}[w]$ அல்லது $\mSat{M}{!C}[w]$} (அல்லது இரண்டும்).}{}
  \tagitem{prvIf}{\indcase{!A}{(!B \lif !C)}{$\mSat{M}{\indfrm}[w]$ எனில்,
    எனில் மட்டுமே, $\mSat/{M}{!B}[w]$ அல்லது $\mSat{M}{!C}[w]$}.}{}
  \tagitem{prvIff}{\indcase{!A}{(!B \liff !C)}{$\mSat{M}{\indfrm}[w]$ எனில்,
    எனில் மட்டுமே, $\mSat{M}{!B}[w]$, $\mSat{M}{!C}[w]$ இரண்டும் மெய்,
    அல்லது $\mSat{M}{!B}[w]$, $\mSat{M}{!C}[w]$ இரண்டுமே மெய்யல்ல}.}{}
  \item\ollabel{defn:sub:mmodels-box}
    \indcase{!A}{\Knows_a !B}{$\mSat{M}{\indfrm}[w]$ எனில், எனில் மட்டுமே,
    $R_a ww'$ ஆகும் எல்லா $w' \in W$ க்கும் $\mSat{M}{!B}[w']$}
  \end{enumerate}
\end{defn}

ஆனால் இங்குதான் நமது அணுகுறவுகளின் மீதான கட்டுப்பாடுகளைப் பற்றி நாம்
சிந்திக்க வேண்டும். ஏனெனில் கூறு~\olref{defn:sub:mmodels-box} இன்படி,
$R_a ww'$ ஆகும் எந்த~$w'$ உம் இல்லாதபோதெல்லாம் !!a{formula}~$\Knows_a !B$
ஆனது~$w$ இல் மெய். இயல்பான வாய்ப்புநிலைத் தருக்கத்தின் அதே கூறே இது; ஓர்
உலகிற்கு அடுத்துறுப்புகள் இல்லாதபோது, எல்லா~$\Box$-வாய்பாடுகளும் அங்கு
வெறுமையாக மெய்யாகின்றன. $\Knows$ ஆனது \emph{அறிவைக்} குறிக்கிறது என்று
நாம் கருதினால், இது உள்ளுணர்வுக்கு மிகவும் முரணானதாகத் தோன்றுகிறது. ஏனெனில்,
ஒரு முகவர் $!A$, $\lnot !A$ இரண்டையும் ஒரே நேரத்தில் அறியக்கூடிய சூழ்நிலை
\emph{எதுவும் இல்லை} என்றே நாம் பொதுவாகக் கருதுகிறோம்.

அறிவுநிலைத் தருக்கத்தில் நமது அணுகுறவு எப்போதும் \emph{தற்சுட்டுப்
பண்புடையதாக} இருப்பதை உறுதிசெய்வது ஒரு தீர்வு. மெய்யான உலகம் ஒரு முகவரின்
தகவலுடன் ஒத்திசைகிறது என்ற கருத்துக்கு இது தோராயமாகப் பொருந்துகிறது.
உண்மையில், அறிவுநிலைத் தருக்கங்கள் பொதுவாக S5 ஐப் பயன்படுத்துகின்றன; ஆனால்
$\Knows_a$ தொடர்பு துல்லியமாக எதைக் குறிக்க வேண்டும் என்பதைப் பொறுத்து,
மற்றவர்கள் வலுக்குறைந்த அமைப்புகளைப் பயன்படுத்தலாம்.

\begin{figure}
  \begin{center}
    \begin{tikzpicture}[modal]
      \node[world] (w1) [label={[align=right]right:\mTrue{p}\\ \mFalse{q}}]{$w_1$};
      \node[world] (w2) [label={[align=right]right:\mTrue{p}\\ \mTrue{q}},
        above right=of w1]{$w_2$};
      \node[world] (w3) [label={[align=right]right:\mFalse{p}\\ \mFalse{q}},
        below right=of w1] {$w_3$};
      \draw[<->] (w1) to node {$a$} (w2);
      \draw[->,loop left] (w1) to node {$a,b$} (w1);
      \draw[<->] (w1) to [swap] node {$b$} (w3);
      \draw[->,loop above] (w2) to node {$a,b$} (w2) ;
      \draw[->,loop below] (w3) to node {$a,b$} (w3) ;
    \end{tikzpicture}
  \end{center}
  \caption{ஓர் எளிய அறிவுநிலை மாதிரி.}
  \ollabel{fig:simple}
\end{figure}

\begin{prob}
  பின்வருவனவற்றில் எவை \olref[aml][el][trw]{fig:simple} இல் மெய்யாகின்றன
  என்பதைக் கருதுக:
  \begin{enumerate}
    \item $\mSat{M}{\lnot q}[w_1]$;
    \item $\mSat{M}{\Knows_a \lnot q}[w_1]$;
    \item $\mSat{M}{\Knows_b \lnot q}[w_1]$;
    \item $\mSat{M}{\Knows_b q \lor \Knows_b \lnot q}[w_2]$;
    \item $\mSat{M}{\Knows_a( \Knows_b q \lor \Knows_b \lnot q)}[w_2]$;
    \item $\mSat{M}{\EKnows_{\{a,b\}} \lnot q}[w_3]$;
    \end{enumerate}
\end{prob}

ஓர் உலகில் மெய்மை என்ற அடிப்படை வரையறையை இப்போது அளித்துவிட்டதால்,
வாய்ப்புநிலைச் செல்லுபடிநிலை, பின்விளைவு போன்ற இயல்பான வாய்ப்புநிலைத்
தருக்கத்தின் பிற பொருண்மையியல் கருத்துகள், வாய்ப்புநிலைச் செயற்குறிகளின்
பொருள்கோளைப் பற்றிய இப்புதிய சிந்தனைமுறைக்குப் பயன்படுத்தப்பட்டபடியே
எளிதாகத் தொடர்கின்றன.

பொதுவறிவுச் செயற்குறி~$\CKnows_G$ க்கான மெய்மை நிபந்தனைகளையும் இப்போது
அளிக்க முடியும். தொடர்பு~$R$ இன் \emph{கடப்பு அடைவு}~$R^+$ ஆனது
\olref[sfr][rel][ops]{sec} இன்படி பின்வருமாறு வரையறுக்கப்படுகிறது என்பதை
நினைவுகூர்க:
\begin{align*}
  R^+ &= \bigcup_{ n \in \mathbb{N}} R^n,
\intertext{இங்கு}
  R^0 & = R \text{ மற்றும்}\\
  R^{n+1} & = \Setabs{\tuple{x, z}}{\exists y (R^n xy \land Ryz)}.
\end{align*}
பின்னர், $G$ ஒரு முகவர் குழுவாக இருக்கும் இடத்தில், எல்லா முகவர்களின்
அணுகுறவுகளின் ஒன்றியத்தின் கடப்பு அடைவாக $R_G = ( \bigcup_{b
\in G} R_b )^+$ ஐ வரையறுக்கிறோம்.

\begin{defn}
$G' \subseteq G$ எனில், $R_{G'} w w'$ ஆகும் ஒவ்வொரு $w'$ க்கும்
$\mSat{M}{!A}[w']$ ஆகும்போது, அப்போதும் அப்போது மட்டுமே,
$\mSat{M}{\CKnows_{G'} !A}[ w]$ எனக் கொள்கிறோம்.
\end{defn}

\end{document}
